\chapter{Preliminarii}

\section{RAG (Retrieval-Augmented Generation)}

Un model lingvistic mare generează răspunsuri pe baza tiparelor învățate la antrenare, fixate în parametrii săi. Pentru informație absentă din datele de antrenare, modelul fie semnalează că nu o deține, fie generează un răspuns plauzibil, dar nesusținut, cu aceeași certitudine ca unul corect (\emph{halucinație} \cite{ji2023hallucination}). În domenii cu miză ridicată (juridic, medical, administrativ), un astfel de răspuns fals poate fi mai dăunător decât recunoașterea lipsei de informație.

\emph{RAG} (Retrieval-Augmented Generation) \cite{lewis2020retrieval} atenuează acest risc schimbând sursa răspunsului: modelul nu mai generează din cunoașterea parametrică, ci pe baza unor documente regăsite în prealabil. Procesul cuprinde două etape: \emph{regăsirea}, în care sistemul selectează automat fragmentele cele mai relevante pentru întrebare, și \emph{generarea}, în care răspunsul este compus pe baza acelor fragmente, oferite modelului ca \emph{context}. Răspunsul devine astfel ancorat în surse, poate fi însoțit de citații verificabile, iar riscul de halucinație scade \cite{gao2023retrieval}. În plus, RAG separă cunoașterea (documentele) de capacitatea de a genera text, astfel încât actualizarea informației nu necesită reantrenarea modelului.

\section{Reprezentări vectoriale și regăsire}

Abordarea clasică de regăsire este cea \emph{lexicală}, în care un fragment este relevant dacă împărtășește cuvinte cu întrebarea, ponderate după frecvență (familia \emph{BM25} \cite{robertson2009bm25}). Metoda este rapidă, dar compară forme de cuvinte, nu sensuri: o întrebare despre „acte necesare" nu se potrivește cu un text despre „documente obligatorii", deși înțelesul este același (\emph{nepotrivire de vocabular}). În limba română problema se accentuează din cauza flexiunii bogate, care multiplică formele de suprafață ale aceluiași cuvânt.

Soluția este compararea sensului. Pornind de la observația că două cuvinte cu contexte asemănătoare au înțelesuri asemănătoare, cuvintele au fost reprezentate prin vectori \emph{denși} (de exemplu \emph{word2vec} \cite{mikolov2013efficient}); modelele neuronale ulterioare au extins ideea la fraze și fragmente întregi, transformând un text într-un singur vector, numit \emph{embedding}. Regăsirea devine astfel \emph{semantică} (\emph{dense retrieval} \cite{karpukhin2020dense}): întrebarea și fragmentele sunt transformate în embeddings de același model (de exemplu multilingvul \emph{BGE-M3} \cite{chen2024bgem3}), iar relevanța este dată de apropierea vectorilor, măsurată prin \emph{similaritate cosinus}. Pentru a evita compararea cu toți vectorii, regăsirea folosește o indexare aproximativă a vecinilor, de exemplu \emph{HNSW} \cite{malkov2020hnsw}, oferită de baze de date vectoriale precum \emph{Chroma} \cite{chroma}.

\section{Modele lingvistice locale}
\label{sec:cuantizare}

Confidențialitatea datelor impune adesea ca modelul să ruleze local, fără transmiterea acestora către servicii externe; principala constrângere devine atunci memoria, determinată de numărul de \emph{biți} alocați fiecărui parametru. Implicit, un parametru ocupă 32 de biți (\texttt{FP32}), deci un model de 3 miliarde de parametri necesită $\sim$12~GB, prea mult pentru un calculator obișnuit. \emph{Cuantizarea} \cite{frantar2023gptq} reduce numărul de biți per parametru: modelul ocupă mai puțin și se evaluează mai rapid, în schimbul preciziei. Nivelurile uzuale pentru un model 3B sunt:

\begin{itemize}
    \item \texttt{FP16} (16 biți): aproximativ jumătate din memorie ($\sim$6~GB), cu pierdere de precizie practic neglijabilă;
    \item \texttt{Q8} (8 biți întregi): $\sim$3~GB, cu degradare redusă a calității;
    \item \texttt{Q4} (4 biți întregi): $\sim$2~GB, cea mai agresivă comprimare uzuală, cu risc real de degradare a raționamentului.
\end{itemize}

Schemele \emph{k-quants} din \texttt{llama.cpp} (sufixul \texttt{\_K\_M}) nu comprimă uniform, ci păstrează la precizie mai mare parametrii importanți, obținând un compromis mai bun. Astfel de modele se rulează local prin \texttt{llama.cpp} (format \texttt{GGUF}) sau prin interfețe precum \emph{Ollama} \cite{ollama}. Pe lângă cuantizare, modelele locale se diferențiază prin \emph{scară} (numărul de parametri: modelele mai mari raționează mai fidel, dar cer mai multe resurse) și prin \emph{specializare lingvistică}: un model generalist multilingv (Qwen2.5 \cite{yang2024qwen25}) acoperă multe limbi, în timp ce unul antrenat special pe o limbă cu resurse reduse (RoLlama3 \cite{masala2024openllmro}) poate compensa subreprezentarea ei.

\section{Verificarea ieșirii}
\label{sec:contracte}

În contexte reglementate, generarea de text liber nu este suficientă: răspunsul trebuie să poată fi verificat. Gheorghe et al.\ (2026) \cite{gheorghe2026robust} arată empiric că aproximativ 63\% dintre erorile injectate la limitele serviciilor externe rămân \emph{nesemnalate} (modelul produce răspunsuri plauzibile pe date incomplete) și propun verificarea sistematică a ieșirii ca direcție de cercetare. Un strat de \emph{verificare a ieșirii} are două componente: o \emph{schemă structurată} (un format tipizat, de exemplu JSON, impus motorului de inferență, care elimină erorile de formă) și un \emph{set de reguli aplicate după generare}, care verifică proprietăți de conținut, precum existența și validitatea citațiilor și prezența în surse a documentelor menționate. La eșec, generarea poate fi reluată cu feedback explicit, iar la epuizarea încercărilor răspunsul devine un refuz controlat.

Verificarea externă se distinge de cea integrată în model (\emph{Self-RAG} \cite{asai2024selfrag}) și se înrudește cu \emph{guardrails} programabile (NeMo Guardrails \cite{rebedea2023nemo}) și cu cadre de evaluare a RAG (RAGAS \cite{es2024ragas}); spre deosebire de un model-judecător extern, regulile de verificare sunt deterministe, auditabile și rulează integral local.
